El desarrollo se apoya en una separación entre la interfaz lógica del registro cuántico y
el backend responsable del almacenamiento del estado. En TMFQSfullstate esa
separación se materializa en la clase \texttt{QuantumRegister}, que conserva la API de
operaciones cuánticas, y en la interfaz \texttt{IStateBackend}, que abstrae la forma de
inicializar, exportar, consultar y mutar el vector de amplitudes.

Esta decisión permite mantener una interacción uniforme con el estado y, al mismo
tiempo, introducir una realización comprimida que administra particionamiento por
bloques, recuperación bajo demanda, persistencia diferida y reutilización temporal. Así,
el simulador conserva su modelo de ejecución compuerta por compuerta, pero delega en
un backend especializado la manera concreta en que las amplitudes residen y circulan en
memoria.
